\PassOptionsToPackage{final}{graphicx}
\documentclass[xcolor=table,8pt]{beamer}
% \usepackage{etex}

\newif\ifmydraft
%\mydrafttrue
\mydraftfalse
%\usepackage{pgfpages}

%\newcommand{\WIP}{\framesubtitle{Draft version of this slide}}
\newcommand{\WIP}{}


\usepackage{expl3}

% \usepackage{polyglossia}
% \setmainlanguage{english}

\usepackage[english]{babel}

\usepackage{csquotes}

\usepackage{fontawesome}

% \usepackage{iflang}

\usepackage{ifthenx}
\usepackage{xstring}

\usepackage{hhline}

% \newcommand{\st}{\textsuperscript{st}}
% \newcommand{\nd}{\textsuperscript{nd}}
% \newcommand{\rd}{\textsuperscript{rd}}
%\newcommand{\th}{\textsuperscript{th}}


%\usepackage[oldenum,olditem]{paralist}

%\addtolength{\parskip}{0.3\baselineskip}
%\addtolength{\topsep}{-5.7\baselineskip}

%\usepackage[backend=biber,style=authoryear,maxnames=5]{biblatex}
% \newcommand{\mytextcite}[1]{%
%   {\structure{\citeauthor{#1}} \citeyear{#1}}}
% \newcommand{\myfootcite}[1]{%
%   \AtNextCite{\DeclareNameAlias{labelname}{first-last}}%
%   \footnote{\structure{\citeauthor*{#1} (\citeyear{#1})}. \citetitle{#1}}%
% }

\newcommand{\poormancite}[2]{%
  {\small% \structure
    {#1}, #2}%
}

\usepackage{import}

%%% Ces packages sont chargés par beamer, mais les ajouter force auctex à les reconnaître
\usepackage{amsfonts, amssymb, amsthm}
\usepackage{amsmath}

\usepackage{mathtools}

%\usepackage{graphicx}

\usepackage{booktabs}
\usepackage{subfigure}
\usepackage{array}
\usepackage{tabu}
\usepackage[justification=centering]{caption}
\usepackage{collcell}

\usepackage{calc} % Pour widthof

\newcommand{\pbox}[2][l]{%
  %\extrarowheight=2pt
  \begin{tabular}[c]{@{}#1@{}}#2\end{tabular}%
  %\extrarowheight=6pt
}

\newcommand{\pfrac}[2]{\frac{\partial #1}{\partial #2}}


% \rowcolors{1}{white}{structure.fg!20!white}

\taburowcolors[2] 2{white .. color1!10!white}

\tabulinesep=_1.5pt^1.5pt
%\def\arraystretch{1.5}

\makeatletter
\AtBeginDocument{
  \renewcommand{\@thesubfigure}{}
  \addtolength\abovedisplayskip{-0.2\abovedisplayskip}%
  \addtolength\belowdisplayskip{-0.2\belowdisplayskip}%
}
\makeatother

% Beamer list spacing
\addtobeamertemplate{itemize/enumerate body begin}{}{\vspace{-0.5em}}

\newcommand\ds{\displaystyle}

\newcommand{\mycovered}[1]{{\color{black!30}#1}}

\usepackage{pdfpages}

\usepackage{pgfplots}
%\usepackage{tikz}
\usetikzlibrary{fpu,arrows ,arrows.meta,bending,
  calc,matrix,shapes.geometric,decorations.pathmorphing,
  intersections,fit,external,backgrounds,
  spy,patterns,tikzmark}

%% Fix pour le souci de bounding box dans pgfplots
\pgfplotsset{compat=1.8}

% \usepackage{forest}
% \forestset{MNRtree/.style={for tree={align=center,child anchor=north, calign=midpoint, anchor=north}}}

\tikzset{
  alerted/.style={color=color2},
  alerted on/.style={alt=#1{alerted}{}},
  alt/.code args={<#1>#2#3}{%
    \alt<#1>{\pgfkeysalso{#2}}{\pgfkeysalso{#3}} % \pgfkeysalso doesn't change the path
  },
}
% \forestset{
%   alerted on/.style={
%     for tree={
%       /tikz/alerted on={#1},
%       edge={/tikz/alerted on={#1}}}}}

\usepackage{xspace}
\usepackage{multirow}
%\usepackage[lined,noend]{algorithm2e}
% \SetArgSty{textsf} % Style pour le texte dans les algorithmes
% \SetKw{KwNot}{not}

\usepackage{comment}

\usepackage{ifdraft}
%\usepackage{datetime2}
\usepackage{siunitx}



% Traductions

\newcommand{\ifnonempty}[3]{% Teste si le premier argument est vide, si oui renvoie le 3e, sinon le 2e
  \def\tempa{}%
  \def\tempb{#1}%
  \ifx\tempa\tempb % Cas vide
  #3 
  \else % Cas non vide
  #2
  \fi}

\newcommand{\texttodo}{\textcolor{red}{TODO}}

%\newcommand{\lang}{english} % Language choice
% \newcommand{\lang}{french} % Language choice
% \newcommand{\translation}[2]{\IfLanguageName{#1}{#2}{}}
% \newcommand{\en}[1]{\translation{english}{#1}}
% \newcommand{\fr}[1]{\translation{french}{#1}}
% \newcommand{\m}[2]{\translation{english}{\ifnonempty{#1}{#1}{\texttodo}}\translation{french}{\ifnonempty{#2}{#2}{\texttodo}}}

%\expandafter\selectlanguage\expandafter{\lang}



% Command definitions


%\setbeamercolor{alerted text}{fg=myred}


\newcommand{\bigO}{O}

\newcommand{\CC}{\mathbb{C}} 
\newcommand{\RR}{\mathbb{R}}
\newcommand{\QQ}{\mathbb{Q}} 
\newcommand{\NN}{\mathbb{N}}
\newcommand{\ZZ}{\mathbb{Z}} 
\newcommand{\KK}{\mathbb{K}}
\newcommand{\FF}{\mathbb{F}} % \F est déjà utilisé pour F5

\DeclareMathOperator{\im}{Im}
\DeclareMathOperator{\Ker}{Ker}         % Noyau

\newcommand{\divides}{\mid}
\newcommand{\setof}[1]{\left\{ #1 \right\}}

%\renewcommand{\tilde}[1]{\widetilde{#1}}
\let\tildeold\tilde
\renewcommand{\tilde}[1]{\skew{3}{\tildeold}{#1}}
\renewcommand{\epsilon}{\varepsilon}
\renewcommand{\phi}{\varphi}
\renewcommand{\emptyset}{\varnothing}
\renewcommand{\setminus}{\smallsetminus}

\newcommand{\grevlex}{\textsc{GRevLex}\xspace}
\newcommand{\lex}{\textsc{Lex}\xspace}
\newcommand{\Wgrevlex}{$W$-\grevlex}
\newcommand{\elim}{\textsc{Elim}\xspace}

\newcommand{\ltgrevlex}{<_\mathsf{grevlex}}
\newcommand{\ltWgrevlex}{<_\mathsf{$W$-grevlex}}

\newcommand{\LT}{\mathsf{LT}}
\newcommand{\LM}{\mathsf{LM}}
\newcommand{\LC}{\mathsf{LC}}
\newcommand{\Spol}{\mathsf{S\text{-}Pol}}

\newcommand{\F}[1]{\ifmmode\mathsf{F_{#1}}\else F$_{#1}$\fi\xspace}
\newcommand{\FGLM}{\ifmmode\mathsf{FGLM}\else FGLM\fi\xspace}

\newcommand{\algo}[1]{\structure{#1}}

\newcommand{\HS}{\mathsf{HS}}
\newcommand{\HI}{i_{\mathsf{reg}}}
\newcommand{\ireg}{\HI}
%\newcommand{\dreg}{d_{\mathsf{reg}}}
\newcommand{\dreg}{\dmax}
\newcommand{\dmax}{d_{\mathsf{max}}}
%\newcommand{\dmaxW}{d_{W,\mathsf{max}}}
%\newcommand{\dmax}{\dreg}
\newcommand{\dmaxW}{d_{W,\mathsf{reg}}}

\newcommand{\sign}{\mathsf{sign}}

\newcommand{\Fext}{F_{\mathsf{ext}}}
\newcommand{\Wdeg}{\deg_{W}}

\newcommand{\pgcdfam}[1]{\mathsf{gcd}\left(#1\right)}

\newcommand{\with}{\text{ with }}

\newcommand{\ordre}[1]{<_{#1}}

\newcommand{\append}{\,\cup} % Ajout d'une ligne à une matrice

\newcommand{\lcm}{\mathsf{lcm}}

\newcommand{\bfe}{\mathbf{e}}
\newcommand{\bff}{\mathbf{f}}
\newcommand{\bfg}{\mathbf{g}}
\newcommand{\bfh}{\mathbf{h}}
\newcommand{\bfp}{\mathbf{p}}
\newcommand{\bfq}{\mathbf{q}}
\newcommand{\bfs}{\mathbf{s}}
\newcommand{\bft}{\mathbf{t}}
\newcommand{\bfz}{\mathbf{z}}
\newcommand{\bfT}{\mathbf{T}}
\newcommand{\bfeps}{\mathbf{\varepsilon}}

\newcommand{\Mon}{\mathsf{Mon}}
\newcommand{\Syz}{\mathsf{Syz}}

\newcommand{\II}{\mathcal{I}}
\newcommand{\GG}{\mathcal{G}}


\renewcommand{\hom}[2][]{%
  \hspace{0.1em}%
  \mathsf{hom}_{#2}^{\ifnonempty{#1}{#1}{\phantom{-1}}}
  \ifnonempty{#1}{\hspace{-0.1em}}{\hspace{-0.2em}}%
}

\newcommand{\truncseries}[1]{\left\lfloor #1\right\rfloor_{\raisebox{0.3ex}{+}}}

\newcommand{\D}{\mathbf{D}}

\newcommand{\merge}[1]{\multirow{2}{*}{#1}}
\newcommand{\QH}{\text{QH}}
\newcommand{\Homo}{\text{$\hom{W}$}}


\newcommand{\lt}{\LT}
\newcommand{\lm}{\LM}
\newcommand{\lc}{\LC}
\newcommand{\sig}{\textsf{sig}}

\newcommand{\lcmlm}{\mathsf{lcmlm}}
\newcommand{\gcdlc}{\mathsf{gcdlc}}
\newcommand{\lcmlt}{\mathsf{lcmlt}}

\newcommand{\Ocomp}[1]{\ensuremath{\operatorname{O}\left(#1\right)}}

\newcommand\paramcolor[1]{{\color{teal}#1}}
\newcommand\varcolor[1]{{\color{violet}#1}}

\newcommand{\paramg}[1]{\paramcolor{\gamma_{#1}}}
\newcommand{\paramG}[1]{\paramcolor{\Gamma_{#1}}}
\newcommand{\paramgg}[1]{\paramcolor{g_{#1}}}
\newcommand{\paramGG}[1]{\paramcolor{G_{#1}}}
\newcommand{\vary}[1]{\varcolor{y_{#1}}}
\newcommand{\varz}[1]{\varcolor{z_{#1}}}
\newcommand{\varq}[1]{\varcolor{q_{#1}}}
\newcommand{\varp}[1]{p_{#1}}
\newcommand{\varXX}{\varcolor{X}}
\newcommand{\varxx}{\varcolor{x}}

\newcommand{\answer}[2][0.3cm]{\textcolor{mygreen}{\hspace{#1}$\rightarrow$ #2}}
\newcommand{\mycitekey}[1]{\structure{#1}}
\newcommand{\strongstructure}[1]{\structure{\textbf{#1}}}


\newcommand{\fracpart}[2]{\frac{\partial #1}{\partial #2}}
%\newcommand{\rk}{\mathrm{rank}}
\DeclareMathOperator{\rk}{rank}
%\DeclareMathOperator{\det}{det}
\DeclareMathOperator{\disc}{disc}

\newcommand{\abs}[1]{\left|#1\right|}

\newcommand{\Sing}{\mathsf{Sing}}
\newcommand{\Crit}{\mathsf{Crit}}
\newcommand{\Jac}{\mathsf{Jac}}
\newcommand{\critval}{\mathsf{CritVal}}

\newcommand{\SingCrit}[2]{K(#1,#2)}
%\newcommand{\SingCrit}[2]{\Sing(#2) \cup \Crit(#1,#2)}
\newcommand{\VarV}{\mathcal{V}}
\newcommand{\VarB}{\mathcal{B}}

\tikzset{annot/.style={draw=darkgreen,thick,rounded corners,anchor=base}}
\newcommand{\ilds}[1]{$\displaystyle #1$}
\colorlet{contourplotscolor}{black}

% \newcommand{\itemdot}{{\usebeamercolor{itemize item}\color{fg}{\usebeamertemplate{itemize item}{}}}\xspace}

\newcommand{\itemdot}{
  \leavevmode\usebeamerfont*{itemize subitem}\raisebox{1.25pt}{\scalebox{0.8}
    {\hbox{\textcolor{color1}{\donotcoloroutermaths$\blacktriangleright$}}}\hspace{0.8ex}}
  % \tikz[baseline=(node.base)]{
  % \node[inner sep=0pt,text height=1.2ex,text width={},circle] (node) {} ;
  % \fill[color1] (node.-120) -- (node.east) -- (node.120) -- cycle; 
  % }
  %   \hspace{0.8ex}
}


\newcommand{\XX}{\,\mathbf{X}}

% For debug:
% \makeatletter
% \newcommand*{\overlaynumber}{\number\beamer@slideinframe}
% \makeatother

% Removed for metropolis
% \setbeamertemplate{itemize item}{\raisebox{1.25pt}{\scalebox{0.8}{$\blacktriangleright$}}}   % Third level item
% \setbeamertemplate{itemize subitem}{\raisebox{1.25pt}{\scalebox{0.8}{$\blacktriangleright$}}}   % Third level item
% \setbeamertemplate{itemize subsubitem}{\raisebox{1.25pt}{\scalebox{0.8}{$\blacktriangleright$}}}   % Third level item
% \setbeamertemplace{itemize item}{\itemdot}

% \makeatletter
% \renewcommand{\tikz@align@newline@}{%  The default definition had 0pt here VVV
%   \unskip \pgfutil@ifnextchar [\tikz@@align@newline {\tikz@@align@newline [5pt]}%
% }
% % For nodes with text width
% \renewcommand{\@xcentercr}[1][5pt]{%
%   \addvspace{-\parskip}\vskip#1\ignorespaces
% }
% \makeatother

\tikzset{
  blocknode/.style={thick,  
    text width=0.95\linewidth},
  defnode/.style={blocknode,draw=color3 ,fill=color3!5!white,
  },
  thmnode/.style={blocknode,draw=color2, fill=color2!5!white}
}

\renewenvironment{theorem}[1][]{%
  \addtobeamertemplate{itemize/enumerate body begin}{}{\vspace{0.5em}}
  \begin{center}
  \begin{tikzpicture}
    \node[thmnode] (thm) \bgroup
      \textcolor{color2}{\textbf{Theorem} #1}\\
    }{%
      \egroup;
      %\draw[thick,color2] (thm.north west) -- (thm.south west);
    \end{tikzpicture}
  \end{center}
}

\renewenvironment{definition}[1][Definition]{%
  \addtobeamertemplate{itemize/enumerate body begin}{}{\vspace{0.5em}}
  %\setbeamercolor{alerted text}{fg=color3}
  \begin{center}
  \begin{tikzpicture}
    \node[defnode] (def) \bgroup
      \textcolor{color3}{\textbf{#1}}\\
    }{
      \egroup;
      % \draw[thick,color2] (def.south west) -- (def.north west) % --++ (2,0)
      % ;
    \end{tikzpicture}
  \end{center}
}


%%% Local Variables:
%%% mode: latex
%%% TeX-master: t
%%% End:


%%% Local Variables:
%%% mode: latex
%%% TeX-master: t
%%% End:
